\clearpage
\fancyhead[L]{\small ASTRONOMICAL YEARBOOK.}
\hypertarget{native-616-lower}{}
\begin{center}\small Göttingische gelehrte Anzeigen. No. 80. Pages $793\;\therefore\;800$. 20 May 1819.\\[3pt]\rule{.72\linewidth}{.3pt}\end{center}

\emph{Astronomical Yearbook for the Year {\upshape 1821}, together with a collection of the latest papers, observations, and reports pertaining to the astronomical sciences. Calculated and published by} Dr. J. E. \textsc{Bode}, \emph{Royal Astronomer, etc.} Berlin 1818. Available from the author and on commission through Ferd. Dümmler. 250 pages in octavo, with one copperplate.

\begingroup\brokenpenalty=10000\relax
Since the arrangement of the astronomical calendar has remained unchanged, we need only review the accompanying papers. The first, by Prof. \textsc{Degen} in Copenhagen, deals with \textsc{Kepler}'s problem and develops the coefficients in the series for the equation of the center two orders further than the paper on the same problem by \textsc{Schubert} in the preceding volume. We refer here to the remarks we made in our notice of the latter. --- Observed eclipses of satellites, occultations of stars, and planetary oppositions by Prof. \textsc{Bürg} at the Vienna Observatory. --- Observed solstices in the years 1812---1817 at the Turin Observatory by \textsc{Plana}. These observations, too, made with a \textsc{Reichenbach} repeating circle with an upright column, 18 inches in diameter, give an obliquity of the ecliptic from the winter solstices that is 4 seconds of arc smaller than that from the summer solstices. When the editor says in a note that this highly remarkable phenomenon has already been investigated and explained by astronomers and physicists, and asserts still more definitely on p. 197 that \textsc{Bessel} eliminated it \emph{by means of an improved refraction table}, he is mistaken, for applying \textsc{Bessel}'s refraction table by no means removes the difference. If anything, the difference turns out to be a little greater than when other refraction tables are used. --- Observed planetary oppositions at Kremsmünster by \textsc{Derflinger}. --- On various astronomical \SourceLeaf{617}subjects by \textsc{Littrow}. For brevity we confine ourselves here to a few remarks on the results that Mr. \textsc{Littrow} drew from the Königsberg observations. Reprinted here is the list, previously published in the Zeitschrift für Astronomie, of the polar distances of 23 fixed stars that Mr. \textsc{Littrow} derived from the Königsberg observations. As is well known, these polar distances are all several seconds of arc greater than those that \textsc{Piazzi}, \textsc{Oriani}, and \textsc{Pond} derived in close agreement from their observations. It is to be hoped that the source of this striking difference will in time be discovered and made to yield lessons for practical astronomy. Yet however much we recognize the exceptional care with which the excellent Königsberg astronomer determined the division errors of his circle, Mr. \textsc{Littrow}'s judgment that the Königsberg observations, above all others, are to be declared entirely free of all constant errors seems to us too hasty and venturesome. Mr. \textsc{Littrow} further presents here the differences in right ascension that he derived from the first two annual volumes of the Königsberg observations for the 36 fundamental stars of \textsc{Maskelyne}, together with the weights and probable errors to be assigned to these determinations. If these determinations really had the reliability attributed to them here, there is no doubt that they would far surpass anything we have hitherto had of this kind. Yet it is precisely this unprecedented accuracy, which with the means hitherto available one could otherwise achieve only through observations spanning a generation and which is claimed to have been attained from two years of observations, that makes a more rigorous examination all the more necessary, since this subject is of the utmost importance for practical astronomy. The latest issue of the Zeitschrift für Astronomie, in which Mr. \textsc{Littrow} gives a fuller account of his procedure, provides the means for this examination. We see from it that Mr. \textsc{Littrow}'s treatment of the observations is itself at variance with the spirit of the method of least squares and, what is much more important, that his determination of the weight of the results runs entirely counter to that spirit in several essential respects and is therefore illusory. Mr. \textsc{Littrow} sets (Zeitschrift für Astronomie, vol. 6, p. 23) the weight of the determination of the difference in right ascension between two stars, from their differences in right ascension with a third, when the latter determinations have weights $p$ and $q$, $=\dfrac{4pq}{p+q}$, which is incorrect. The correct formula is only $\dfrac{pq}{p+q}$. This alone would already reduce the weights of his final results to nearly a quarter of their purported value. Yet still more consequential is the circumstance that Mr. \textsc{Littrow} overlooked, that the combination of the partial results into a final result and the determination of the weight of the latter are valid only on the assumption that the former must all be entirely independent of one another, which is by no means the case here. This reduces Mr. \textsc{Littrow}'s large numbers so much that, e.\,g., the determination of the difference in right ascension between $\gamma$ in Pegasus and $\alpha$ in Cygnus, which is claimed to have the weight of 1507 individual comparisons, would scarcely retain much more than the weight of 40 individual comparisons. Some other aspects of Mr. \textsc{Littrow}'s calculation might also give rise to objections, which we must pass over here for brevity. --- Opposition of Pallas in 1816, and other astronomical observations in 1817 by \textsc{David} and \textsc{Bittner} at the Prague Observatory. We give here only a few remarks on the zenith distances of the Sun and of fixed stars observed with a 12-inch \textsc{Reichenbach} repeating circle. Observations of this kind also appeared in several earlier volumes of the Yearbook. It would have been desirable to present these observations in a more tabular form, which would make it easier to see what conclusions might be drawn from them. Mr. \textsc{David} is in the habit of comparing the observed zenith distance with the one he calculates from the altitude of the pole, assumed known, and the celestial body's declination \SourceLeaf{618}taken from star catalogues or the solar tables. He regards the difference as the observed refraction and the departure of that observed refraction from the tabulated refraction as a correction to the latter. We cannot approve of this procedure. As far as we know, Mr. \textsc{David} determined the altitude of the pole at Prague by \textsc{Hell}'s method, using stars that culminate at nearly the same altitude to the north and south and whose declinations he takes from the star catalogues. The method is good, but only for observatories and instruments of lesser rank. Anyone undertaking so extremely delicate a task as improving the refraction tables must first establish the altitude of the pole on their own and independently of determinations by others. What Mr. \textsc{David} regards as errors in the refraction tables arose from the combined effect of the error in the altitude of the pole, errors in the declinations, errors in the instrument, errors in the observations, errors in the reductions for aberration and in the other reductions, and perhaps, to a very small extent, errors in the refraction tables. A few isolated observations are insufficient to separate these errors from one another. What is required is a series of observations carried out systematically with the best instruments over a period of years and put to systematic use. Mr. \textsc{David} believes that even if one were to doubt the absolute accuracy of the distances themselves, the influence of refractions that vary with changes in temperature would still be apparent in the observed zenith distances in relation to one another. Yet these changes are already taken into account when the tables themselves are applied. It has been proved theoretically that, however the density of the layers of air may change, at such moderate zenith distances refraction depends solely on the density of the lowest layer of air, provided that all the layers are horizontal. There is no doubt that this assumption is not always strictly true. Instead of a refraction that follows changes in temperature, what results \emph{from that} is an irregular refraction that defies calculation. Its adverse effects can be overcome only by making frequent observations and by excluding those in which the celestial bodies gave particularly unsteady images. We therefore see in the differences mentioned only the workings of the unavoidable irregularities in the observations. It also struck us that Mr. \textsc{David} always reduces the zenith distances observed with the crystal cube (prism) by $1''$. We cannot approve of this purported correction, since we find nothing in the nature of the instrument to justify assuming a constant error for all zenith distances. --- Also from the same observatory, observations of Jupiter, Uranus, Saturn, and Ceres. --- Comet of 1 November 1817, discovered by Dr. \textsc{Olbers} but observed only once. --- Observations of the (first) comet of 1818 by Dr. \textsc{Olbers}. --- On some remarkable regions of the Milky Way, by Dr. \textsc{Herschel} (excerpt from the Philosoph. Transact.). --- Observations of planetary oppositions, etc., in Vilna by \textsc{Sniadecki}. --- Observations and determinations of the orbit of the first comet of 1818, by \textsc{Enke}. What is remarkable about this comet is its faintness, which increased when, according to the calculation, its light (whether regarded as its own or as borrowed) should still have been increasing. The comet probably underwent considerable internal changes during the period when it was visible. --- Observed eclipses of satellites, occultations of stars, and sunspots in 1817, by \textsc{Hallashka} in Prague. --- Formulas by \textsc{Bessel} for calculating the position of a celestial body from observed alignments with four stars. The formulas are based essentially on first determining where the two great circles containing the two known pairs of stars intersect the equator (or the ecliptic), although the solution is given an elegant, purely analytical form. --- Observations by the same astronomer of planetary oppositions, stellar occultations, solstices, and the Pole Star in 1817. --- Extract from the astronomical observations at the Berlin Observatory in the same year. --- Determination of the altitude of the celestial pole at Kraków. With \SourceLeaf{619}a repeating theodolite by \textsc{Reichenbach}, Mr. \textsc{Lorenz}, a district commissioner, found $50^\circ\,3'\,57''7$ from observations of the North Star and $50^\circ\,3'\,50''$ from observations of the Sun. Mr. \textsc{Sniadecky}'s earlier determination agrees very closely with the mean. Later, \textsc{Littrow} reduced that determination by $17''$, apparently without justification, judging by the new observations. --- On the accuracy of astronomical observations and calculations. This essay aims to show how very much harder it is nowadays to make new advances in astronomy than it was formerly, when so much less was demanded of the astronomer in observing and calculating. One might envy the position of the earlier astronomers, says the author. They could observe and calculate the motions of the heavens with far less effort, since no great accuracy, let alone the greatest accuracy, could yet be expected, and they nevertheless won fame among their contemporaries and with posterity, and so on. True as this is when we consider the achievements of earlier astronomers from \emph{our} point of view, we must also bear in mind how unwieldy and imperfect their instruments were, how laborious their methods of observation were, and how little their methods of calculation had developed. Giving this due consideration, one may well doubt that \textsc{Tycho}, for example, attained his well-deserved renown as a practical astronomer any more easily than \textsc{Piazzi}. One could still raise a number of objections to particular remarks in this essay about present-day practical astronomy. After calculating, for example, that on a circle with a radius of 12 inches a mark one hundredth of a line thick already covers more than 14 arcseconds, the author asks, “Could a human hand still manage to divide a duodecimal line into 100 distinguishable marks or dots?” But there is no requirement actually to divide the decimal line with 100 marks. The foremost instrument makers can make the divisions in whatever number is suitable, with such accuracy that the errors amount to \emph{much} less than one hundredth of a line, as the masterpieces of \textsc{Reichenbach}, \textsc{Repsold}, and \textsc{Troughton} demonstrate \emph{in practice}. --- New elements of Juno's orbit, observations of this planet and of the Sun, and other astronomical reports by \textsc{Nicolai}. The reports mention the \textsc{Repsold} circle at the observatory here, but the diameter of this instrument, given as 5 feet perhaps through a printing error, should be corrected to $3\tfrac{1}{2}$ Paris feet. --- Observation of the lunar eclipse of 20 April 1818 in Dresden. --- Right ascensions of 35 of the principal stars for 1818, according to \textsc{Maskelyne}'s last observations, from the Nautical Almanac for 1820. (We note here that the Nautical Almanac for 1821 contains a still more extensive catalogue based on \textsc{Pond}'s observations with the new transit telescope.) Observations of the apparent magnitude of various stars from the years 1704---1709 (from the papers left by the Berlin astronomer \textsc{Gottfr. Kirch}). --- On the \textsc{Repsold} meridian circle, together with various astronomical observations by \textsc{Gauss}. --- Observations of the solar eclipse of 5 May 1818 by \textsc{Struve}. --- Observation of the opposition of Vesta in 1818 by \textsc{Enke}. --- Further, ephemerides for Vesta, Juno, and Pallas by \textsc{Enke}, \textsc{Posselt}, \textsc{Nicolai}, and \textsc{Dirksen}. --- Space does not permit us to list the remaining shorter articles and notes individually here.
\par\clearpage\endgroup
